\section{Community Engagement}

\subsection{Community Relations Framework}

The company has established a comprehensive community relations framework that aligns with Indonesian regulatory requirements and international best practices. This framework encompasses multiple stakeholder groups including local communities, indigenous peoples, and vulnerable populations affected by company operations. The framework is structured around four key pillars: communication, consultation, grievance mechanisms, and benefit-sharing arrangements.

The company's community engagement strategy is guided by the principle of Free, Prior, and Informed Consent (FPIC), particularly when operating in areas with indigenous populations. This approach ensures that communities have the right to negotiate the terms under which they would be willing to accommodate projects that may affect their land, resources, or way of life. The FPIC process is conducted in accordance with the Ministry of Environment and Forestry Regulation No. P.14/MENLHK/SETJEN/KUM.1/3/2017 concerning Social Forestry Management.

\subsection{Community Development Programs}

The company implements various community development programs across its operational areas in Indonesia. These programs are designed to address local needs and priorities while promoting sustainable development. Key initiatives include:

Education and Skills Development: The company has established vocational training centers in partnership with local technical schools, providing skills training in areas such as welding, heavy equipment operation, and hospitality management. These programs have benefited over 500 local residents since 2018, with many graduates securing employment within the company or its contractors.

Healthcare Access: Mobile health clinics are deployed to remote communities, offering basic medical services, maternal and child health care, and health education. The company has also constructed and equipped three community health centers in collaboration with the Ministry of Health, improving access to healthcare for approximately 15,000 residents.

Economic Empowerment: Small and medium enterprise (SME) development programs provide business training, micro-financing, and market access support to local entrepreneurs. The company has facilitated the establishment of cooperative businesses in agriculture, fisheries, and handicrafts, creating sustainable income opportunities for community members.

Infrastructure Development: Community infrastructure projects include the construction and maintenance of roads, bridges, and clean water facilities. These projects are implemented in consultation with local governments and communities to ensure they address priority needs and are sustainable in the long term.

\subsection{Stakeholder Engagement Mechanisms}

The company employs multiple channels for ongoing stakeholder engagement, ensuring that community voices are heard and incorporated into decision-making processes. These mechanisms include:

Community Advisory Panels: Established in each operational area, these panels consist of community leaders, local government representatives, and civil society organizations. They meet quarterly to discuss project developments, address concerns, and provide input on community development initiatives.

Community Information Centers: Physical information centers are maintained in each operational area, providing accessible information about company activities, environmental monitoring results, and grievance procedures. These centers are staffed by community members trained in information management and customer service.

Digital Engagement Platforms: Recognizing the increasing importance of digital communication, the company has developed online platforms for community engagement, including mobile applications for reporting concerns and accessing information. These platforms are designed to be accessible to users with varying levels of digital literacy.

Annual Stakeholder Forums: Large-scale forums are held annually, bringing together community representatives, government officials, NGOs, and company management to discuss performance, challenges, and opportunities for collaboration. These forums provide a transparent platform for dialogue and accountability.

\subsection{Grievance Management System}

The company has implemented a robust grievance management system that complies with the Indonesian National Standards for Complaint Handling (SNI 8000:2015). The system includes multiple reporting channels, including toll-free hotlines, email, in-person meetings, and community liaison officers. All grievances are logged, tracked, and resolved within agreed timeframes, with regular reporting to relevant stakeholders.

The grievance mechanism is designed to be accessible to all community members, including those with disabilities, women, and marginalized groups. Special provisions are made for non-literate stakeholders, including oral grievance submission and pictorial information materials. The company has also established partnerships with local NGOs to provide independent grievance handling support for vulnerable community members.

\subsection{Benefit-Sharing Arrangements}

The company has developed benefit-sharing arrangements that ensure communities receive tangible benefits from company operations. These arrangements include:

Community Development Agreements: Formal agreements with local communities outline specific commitments for infrastructure development, employment opportunities, and business development support. These agreements are negotiated transparently and include dispute resolution mechanisms.

Local Content Policies: The company prioritizes local employment, procurement, and contracting wherever possible. Local content targets are set for each operational area, with regular reporting on progress and challenges. The company also provides training and capacity building to enhance local competitiveness.

Revenue Sharing: In accordance with Indonesian regulations, the company contributes to regional and village funds based on production volumes. These funds are managed by local governments and communities for development projects that benefit the broader population.

Environmental Compensation: Where company operations impact natural resources, the company provides environmental compensation through reforestation programs, biodiversity conservation initiatives, and sustainable resource management projects. These programs are developed in consultation with environmental experts and local communities.

\subsection{Performance Monitoring and Reporting}

The company conducts regular assessments of its community engagement performance using both quantitative and qualitative indicators. These assessments include:

Community Satisfaction Surveys: Annual surveys measure community perceptions of company performance, identifying areas of strength and opportunities for improvement. Survey results are disaggregated by demographic groups to ensure all voices are represented.

Third-Party Audits: Independent auditors assess the company's compliance with community engagement commitments and international standards. Audit findings are published in sustainability reports and shared with relevant stakeholders.

Impact Evaluations: Comprehensive evaluations are conducted periodically to assess the long-term impacts of community development programs on livelihoods, social cohesion, and environmental sustainability. These evaluations inform program adjustments and strategic planning.

The company publishes an annual Community Engagement Report that provides detailed information on performance against targets, grievance statistics, benefit-sharing outcomes, and case studies of successful community partnerships. This report is made publicly available and shared with government agencies, community representatives, and other stakeholders.

\subsection{Challenges and Continuous Improvement}

Despite significant progress, the company recognizes ongoing challenges in community engagement, including:

Cultural Sensitivity: Operating across diverse cultural contexts in Indonesia requires continuous learning and adaptation of engagement approaches. The company invests in cultural competency training for staff and partners with anthropologists and cultural experts to improve understanding.

Land Rights Conflicts: Complex land tenure systems and historical claims can lead to conflicts over resource use. The company has established a dedicated land rights team to address these issues through mediation, legal support, and alternative dispute resolution mechanisms.

Economic Disparities: Addressing economic disparities between company operations and surrounding communities remains a challenge. The company is exploring innovative approaches to inclusive economic development, including impact investing and social enterprise models.

The company is committed to continuous improvement in community engagement, regularly reviewing and updating its approaches based on stakeholder feedback, emerging best practices, and changing community needs. This commitment is reflected in the company's sustainability strategy and performance targets.